% =============================================================================
\section{Arquitectura computacional del \textit{pipeline}}
\label{sec:arquitectura}
% =============================================================================

El \textit{pipeline} está implementado en \python{}~3.12 e integra \lammps{}
mediante subprocesos. La Figura~\ref{fig:arquitectura} resume la estructura
modular.

\begin{figure}[ht]
    \centering
    % Sustituir por figura real cuando esté disponible
    \fbox{\parbox{0.85\linewidth}{\centering\vspace{1.5cm}
        \textit{[Diagrama de módulos: config.py $\to$ simulation.py $\to$
        analysis.py $\to$ visualization.py / runner.py]}
        \vspace{1.5cm}}}
    \caption{Arquitectura modular del \textit{pipeline}. Las flechas indican
    dependencias de importación.}
    \label{fig:arquitectura}
\end{figure}

\subsection{Gestión de parámetros: metaclase \texttt{ConfigMeta}}
\label{sec:configmeta}

El problema: cada nueva simulación requiere validar docenas de parámetros
(temperatura, velocidad, ángulo, tamaño de celda...) de forma robusta, y poder
serializar la configuración exacta para reproducibilidad posterior. La solución
ad hoc habitual son diccionarios sin tipo ni validación.

La solución adoptada es una \textbf{metaclase} \verb|ConfigMeta| que genera
automáticamente \verb|__init__|, \verb|__repr__|, \verb|__eq__|,
\verb|to_dict|/\verb|from_dict| y \verb|to_json|/\verb|from_json| para
cualquier clase de configuración que la use como metaclase:

\begin{lstlisting}[caption={Fragmento de \texttt{config.py}: metaclase y descriptor.},
                   label={lst:configmeta}]
class FieldDescriptor:
    """Descriptor de datos con coercicion de tipo y validacion."""
    def __set__(self, obj, value):
        coerced = self.type_(value)
        if self.validator and not self.validator(coerced):
            raise ValidationError(self.name, coerced,
                                  self.constraint_msg)
        obj.__dict__[self.name] = coerced

class ConfigMeta(type):
    """Metaclase: genera __init__, to_json, from_json automaticamente."""
    def __new__(mcs, name, bases, namespace):
        fields = {k: v for k, v in namespace.items()
                  if isinstance(v, FieldDescriptor)}
        namespace['_fields'] = fields
        # Genera __init__ con firma tipada
        namespace['__init__'] = mcs._make_init(fields)
        namespace['to_dict']  = mcs._make_to_dict(fields)
        # ... (from_dict, to_json, from_json)
        return super().__new__(mcs, name, bases, namespace)

class ClusterImpactConfig(metaclass=ConfigMeta):
    N_cluster = FieldDescriptor(int,   validator=lambda x: x>0)
    v0        = FieldDescriptor(float, validator=lambda x: 0<x<=20)
    angle     = FieldDescriptor(float, validator=lambda x: 0<=x<90)
    T_target  = FieldDescriptor(float, validator=lambda x: x>0)
    eps_ps    = FieldDescriptor(float, validator=lambda x: x>0)
\end{lstlisting}

\textbf{Beneficio medible}: cualquier parámetro fuera de rango levanta una
\verb|ValidationError| tipada antes de lanzar \lammps{}, con el mensaje exacto
del parámetro fallido. La reproducibilidad se garantiza guardando el JSON de
la configuración en cada directorio de salida.

\subsection{Jerarquía de excepciones}
\label{sec:excepciones}

Se implementa una jerarquía de cuatro ramas derivadas de
\verb|ReboundMDError|: \verb|IOError|, \verb|SimulationError|,
\verb|ConfigError| y \verb|PhysicsError|. Cada excepción porta un diccionario
\verb|context| con los valores relevantes y redefine \verb|__str__|:

\begin{lstlisting}[caption={Excepción con contexto estructurado.}]
class SimulationRuntimeError(SimulationError):
    def __init__(self, returncode, script, lammps_output):
        self.context = {'returncode': returncode,
                        'script': script}
        self.lammps_tail = lammps_output[-20:]

    def __str__(self):
        tail = '\n  '.join(self.lammps_tail)
        return (f"LAMMPS terminó con código {self.context['returncode']}"
                f" ejecutando '{self.context['script']}'.\n"
                f"Últimas líneas:\n  {tail}")
\end{lstlisting}

\textbf{Beneficio}: en un barrido de 1\,350 simulaciones en paralelo, cuando
falla un run el mensaje de error indica exactamente qué parámetro causó el
fallo y las últimas 20 líneas del log de \lammps{}, sin necesidad de abrir
el archivo de log manualmente.

\subsection{Decoradores: instrumentación sin modificar el código científico}
\label{sec:decoradores}

Cuatro decoradores instrumentan los métodos del analizador sin modificar la
lógica física:

\begin{lstlisting}[caption={Decoradores en \texttt{analysis.py}.}]
@timed          # mide tiempo de ejecucion por metodo
@require_data   # lanza error si no se han cargado datos
@cached         # cachea resultado; se invalida al recargar datos
@validate_physical('COR', 0.0, 1.05)  # fabrica parametrica
def compute_all(self) -> ClusterReboundResult: ...
\end{lstlisting}

El decorador \verb|@cached| es particularmente valioso: en el script de
verificación que carga \num{1350} archivos de datos, cada propiedad derivada
(COR, $\alphaint$, $\We$...) se calcula una sola vez y se reutiliza en las
múltiples figuras, reduciendo el tiempo de análisis en $\approx$70\%.

El decorador \verb|@validate_physical| es una \textbf{fábrica de decoradores}
(decorador paramétrico), que genera un validador específico para cada
observable con sus límites físicos:

\begin{lstlisting}
@validate_physical('alpha_E_int', -0.5, 2.0)
def alpha_E_int(self) -> float:
    return (self.ke_internal_final
            - self.ke_internal_initial) / self.E_incident
\end{lstlisting}

\subsection{Generadores: iteración perezosa sobre trayectorias}
\label{sec:generadores}

Para análisis frame a frame de trayectorias (hasta $\sim$9\,000 pasos por run),
se implementan dos generadores en \verb|ClusterReboundAnalyzer|:

\begin{lstlisting}[caption={Generadores de trayectoria.}]
def iter_frames(self) -> Generator[ClusterFrame, None, None]:
    """Yield ClusterFrame para cada paso registrado."""
    d = self._data
    for i in range(len(d)):
        yield ClusterFrame(time=d.time[i], z_cm=d.z_cm[i],
                           vz_cm=d.vz_cm[i], ...)

def iter_contacts(self) -> Generator[ClusterFrame, None, None]:
    """Yield solo los frames durante la fase de contacto."""
    mask = self._contact_mask()
    for i in np.where(mask)[0]:
        yield ClusterFrame(...)
\end{lstlisting}

\textbf{Uso práctico}: extraer la máxima fuerza de contacto o la duración
del impacto se reduce a una sola línea:
\begin{lstlisting}
t_contact = max(f.time for f in az.iter_contacts())
ke_peak   = max(f.ke_internal for f in az.iter_contacts())
\end{lstlisting}

La evaluación perezosa (lazy) evita materializar el arreglo completo cuando
solo se necesita un estadístico del subconjunto de frames de contacto.

\subsection{Análisis tabular con \texttt{pandas}}
\label{sec:pandas}

El script \verb|verificar_resultados.py| carga los 1\,350 archivos de resultados
en un único \verb|DataFrame| de \texttt{pandas}:

\begin{lstlisting}[caption={Carga masiva y operaciones de agregacion.}]
df = load_all_results()   # -> DataFrame con 13 columnas

# Fraccion de sticking promedio por (T*, eps, N, v0)
stk = (df.groupby(['T','eps','N','v0'])['sticking']
         .mean().reset_index())

# Ley de potencia COR ~ A * v0^{-beta} con scipy
from scipy.optimize import curve_fit
popt, pcov = curve_fit(lambda v, A, b: A*v**(-b),
                       v0_arr, cor_arr, p0=[1.0, 0.5])
beta, sigma_beta = popt[1], np.sqrt(pcov[1,1])
\end{lstlisting}

Frente a la implementación anterior con \verb|list[dict]| y bucles manuales, la
versión \texttt{pandas} reduce el código de agregación de $\sim$80 líneas a
$\sim$15, elimina los errores de indexación y expone directamente los
estadísticos con \verb|.describe()|.

El ajuste con \verb|scipy.optimize.curve_fit| (mínimos cuadrados no-lineales,
internamente basado en LAPACK \texttt{dgelsd}) reemplaza la regresión
log-log linealizada, produciendo estimaciones $\beta \pm \sigma_\beta$
reportables en tablas de resultados.

\subsection{Barrido paramétrico paralelo}
\label{sec:paralelo}

El runner ejecuta el barrido de 1\,350 simulaciones mediante
\verb|ProcessPoolExecutor| con \verb|max_workers=4|:

\begin{lstlisting}[caption={Barrido paralelo en \texttt{runner.py}.}]
with ProcessPoolExecutor(max_workers=jobs) as pool:
    futures = {pool.submit(_run_one, cfg): cfg
               for cfg in all_configs}
    for fut in as_completed(futures):
        cfg = futures[fut]
        try:
            result = fut.result()
            ok += 1
        except SimulationRuntimeError as exc:
            logging.error("Job %s falló: %s", cfg.tag, exc)
            fail += 1
\end{lstlisting}

Al capturar \verb|SimulationRuntimeError| (en lugar del genérico
\verb|Exception|), los fallos individuales no detienen el barrido completo y
quedan registrados con contexto suficiente para diagnóstico posterior sin
perder los runs exitosos.

\subsection{Resumen de patrones aplicados}

La Tabla~\ref{tab:patrones} relaciona cada patrón de programación avanzada
con el problema de investigación que resuelve.

\begin{table}[ht]
    \centering
    \small
    \caption{Patrones de programación avanzada y sus beneficios en el
    \textit{pipeline}.}
    \label{tab:patrones}
    \begin{tabular}{p{3.2cm} p{4.0cm} p{4.5cm}}
        \toprule
        Patrón & Problema resuelto & Beneficio medible \\
        \midrule
        Metaclase \texttt{ConfigMeta} &
            Validar $>20$ parámetros sin código repetitivo &
            Error tipado antes de lanzar \lammps{}; JSON reproducible \\[4pt]
        Jerarquía de excepciones &
            Distinguir errores de IO, LAMMPS, física y config &
            Diagnóstico en barridos paralelos sin abrir logs \\[4pt]
        \texttt{@timed}, \texttt{@cached} &
            Evitar recomputar sobre 1\,350 archivos &
            Reducción $\approx$70\% en tiempo de análisis \\[4pt]
        \texttt{@validate\_physical} &
            Detectar resultados no físicos automáticamente &
            Alerta temprana antes de generar figuras \\[4pt]
        Generadores &
            Análisis frame a frame sin materializar arrays &
            Una línea para estadísticos de fase de contacto \\[4pt]
        \texttt{pandas} &
            Agregación y pivoteo sobre 1\,350 filas &
            De 80 líneas de bucles a 15 líneas declarativas \\[4pt]
        \texttt{scipy.optimize} &
            Ajuste no-lineal COR$\sim Av_0^{-\beta}$ &
            $\beta \pm \sigma_\beta$ reportable; sin sesgo de linealización \\[4pt]
        \texttt{ProcessPoolExecutor} &
            Barrido de 1\,350 simulaciones &
            Factor $\approx$4$\times$ de aceleración en 4 núcleos \\
        \bottomrule
    \end{tabular}
\end{table}
